\section{Angular transfer and its two leading modes}\label{sec:angular}

Throughout this section, $1\le q<4$, $d=\sqrt q=2\cos\lambda$,
$\rho=\lambda/\pi\in(0,1/3]$, $v=\pi/2$, and $Q=-e^{i\lambda}$.
The argument has two parts. We first construct a Hilbert space from
infinite-plane row correlations and prove an angular estimate there. We
then identify its scalar products with reflection across diagonal cuts
of the physical square array. Boundary conditions enter the resulting
estimate through the norms of the exterior states.

\subsection{Rows, twists, and the estimate}

The local row algebra is the six-vertex form of the planar loop
transfer formalism \citep{TemperleyLieb1971,BaxterKellandWu1976}.
Its finite identities are verified below; the infinite-plane
reflection space and the angular decay estimate require the additional
analytic arguments of this section.

In the spin basis $+,-$, set
\[
 J=\begin{pmatrix}0&1\\1&0\end{pmatrix},\qquad
 L_-=|-\rangle\langle+|,\quad L_+=|+\rangle\langle-|,
 \qquad K(\zeta)=\operatorname{diag}(\zeta,\zeta^{-1}).
\]
Here and below twists satisfy $|\zeta|=1$. In the ordered pair basis
$++,+-,-+,--$, let
\begin{equation}\label{bulk:R}
 R(u)=\begin{pmatrix}
 a(u)&0&0&0\\0&b(u)&1&0\\0&1&b(u)&0\\0&0&0&a(u)
 \end{pmatrix},\qquad
 a(u)=\frac{\sin(\lambda-u)}{\sin\lambda},\quad
 b(u)=\frac{\sin u}{\sin\lambda}.
\end{equation}
Time and space steps have unit length and angles $\alpha_t$ and
$\omega_j$. Their successive sums give a rhombic tile array; initially
$0<\omega_j-\alpha_t<\pi$. A spin is positive when its arrow points
along the left normal of the corresponding step. If $h_0,h_1$ are the
two time sides before and after advancing in space, and $s_0,s_1$ the
two space sides before and after advancing in time, the tile tensor
has output indices $(h_0,s_1)$, input indices $(h_1,s_0)$, and value
$R(\rho(\omega_j-\alpha_t))$.

On a spatial circle of even size $N$, write
\begin{equation}\label{bulk:finite-row}
 T_N^D(\alpha)=\operatorname{Tr}_c
 D_cR_{c0}(\rho(\omega_0-\alpha))\cdots
 R_{c,N-1}(\rho(\omega_{N-1}-\alpha)).
\end{equation}
The auxiliary index is $c$. At the seam, $D_{x,y}$ uses the spin $x$
from tile $-1$ and $y$ from tile $0$. Let $\Omega_N$ be the positive
unit Perron vector of the untwisted rows on the zero-total-spin
sector, and $\Lambda_N(\alpha)$ their eigenvalue. Normalized rows are
$A_N^D(\alpha)=T_N^D(\alpha)/\Lambda_N(\alpha)$.
Their common Perron data, and the limiting operators $A^D(\alpha)$,
are constructed below. In the homogeneous transverse array
$\omega_j=v$, put
\begin{equation}\label{bulk:MN}
 \begin{aligned}
 P_\zeta(z)=A^{K(\zeta)}(z),\qquad
 M(z)&=e^{i\tau\rho z}A^D(z),\\
 N(z)&=M(z)P_\zeta(z)+P_\xi(z)M(z),
 \end{aligned}
\end{equation}
where
\begin{equation}\label{bulk:charges}
 \zeta/\xi=Q^\epsilon,\quad \epsilon\in\{-1,1\},\qquad
 (D,\tau)=(L_-,\epsilon)\ \text{or}\ (L_+,-\epsilon).
\end{equation}
The two summands of $N$ describe the two possible positions of a
source in two consecutive rows of the same angle. Products act from
right to left: rows before the source carry twist $\zeta$, and rows
after it carry twist $\xi$. Thus a source in the second row gives
$MP_\zeta$, whereas one in the first gives $P_\xi M$. The phase in
$M$ records the direction of the source row. Section~\ref{sec:current}
will identify this sum with a two-step path increment; here we prove
the row estimate independently of the path's exterior geometry.

\begin{theorem}[Angular transfer estimate]\label{bulk:angular-estimate}
Fix a compact interval $I\subset(-v,v)$. Let $B_{\rm in}$ be a
product of $P_\zeta$ rows and $B_{\rm out}$ a product of $P_\xi$ rows,
at real angles in $[-v,v]$. Suppose each product contains at least
$n\ge1$ angles in $I$. Uniformly for $x,y\in[-v,v]$,
\begin{align}
 \|B_{\rm out}N(x)B_{\rm in}\|&\le Cn^{-1+\rho},
 \label{bulk:low-estimate}\\
 \|B_{\rm out}\{N(x)-e^{i\tau(x-y)}N(y)\}B_{\rm in}\|
 &\le Cn^{-1-\rho}.\label{bulk:high-estimate}
\end{align}
The constant depends on $q$ and $I$. Boundedly many additional pairs
of before-twist rows may be inserted immediately to the right of
$N$, and additional pairs of after-twist rows immediately to its
left. The two terms in \eqref{bulk:high-estimate} may use different
such insertions; on each side their counts must be even. The
constant may depend on a bound for these counts. Multiplication of
$D$ by a common scalar multiplies the bounds by its absolute value.

The same estimates hold in the physical square array sliced in a
diagonal time direction: its time angles are $\pm\gamma_0$, where
$\gamma_0=\pi/4$, and its space angles alternate between
$\pi-\gamma_0$ and $\gamma_0$. The exterior states may contain any
finite number of physical row and spin insertions outside the
sandwich. A scalar contraction is bounded by the displayed rate
times the two reflection norms, with the squares centered on a cut
included in its reflection kernel.
\end{theorem}

The improvement in the second estimate comes from removing one angular
mode, not from assuming that the exterior states have bounded norms.
We first construct the plane reflection space and continue its rows
analytically in the angle. Boundary unitarity then controls the scalar
spectral functions near modulus one. A product of $n$ strict rows
selects spectral energies of order $n^{-1}$. On those low bands,
endpoint reversal singles out the mode $e^{i\tau x}$ with bound
$O(n^{-1+\rho})$; the next odd mode and the remaining errors have
bound $O(n^{-1-\rho})$.
The combination in \eqref{bulk:high-estimate} cancels the former.
The final subsection transfers this estimate, with its exterior norms,
to diagonal cuts of the physical square array.

\subsection{Row algebra and the infinite-plane state}

The finite-dimensional identities used throughout are
\begin{align}
 R_{12}(x-y)R_{13}(x)R_{23}(y)
 &=R_{23}(y)R_{13}(x)R_{12}(x-y),\label{bulk:YB}\\
 R(u)R(-u)&=(1-b(u)^2)I,\label{bulk:inversion}\\
 J_cR(u)^{t_c}J_c&=R(\lambda-u).\label{bulk:crossing}
\end{align}
For completeness, the nontrivial one-minus block of
\eqref{bulk:YB} is symmetric precisely by the three sine identities
\[
 b(x-y)+a(x)b(y)=b(x)a(y),\qquad
 b(y)+a(x)b(x-y)=b(x)a(x-y),
\]
\[
 a(x)+b(x-y)b(y)=a(x-y)a(y).
\]
Spin reversal gives the one-plus block; the remaining blocks are
scalar. The other identities follow directly from \eqref{bulk:R}.
Sliding the intertwiner in \eqref{bulk:YB} through two complete rows
shows that transfers with a common diagonal twist commute: the
intertwiner preserves total spin and hence commutes with the tensor
square of that twist. For arbitrary insertions $D,E$, the same
operation conjugates $D\otimes E$ by a two-spin intertwiner.
Expanding the conjugated matrix in matrix units expresses the
exchanged rows as a finite sum with new insertions. Its coefficients
are analytic wherever that intertwiner is invertible. Division by
the scalar Perron eigenvalues preserves this exchange identity.

For strict angles, the untwisted transfer is primitive on each
total-spin sector. It has positive diagonal entries; moreover an
auxiliary spin can swap with any chosen opposite pair of site
spins, without swapping at the other sites. Such exchanges connect
the sector. Commutation therefore gives common positive left and
right Perron vectors. In the homogeneous array, auxiliary
transposition reverses the product in \eqref{bulk:finite-row}, since
the distinct site operators commute. Equation \eqref{bulk:crossing}
then gives
\begin{equation}\label{bulk:adjoint}
 A^D(\alpha)^*=A^{J\overline D J}(-\alpha)
 \qquad(\alpha\in\mathbb R).
\end{equation}
The corresponding finite-volume identity makes the left and right
vacua equal.

More explicitly, a formal row word $h$ specifies a finite list
$(D_1,\alpha_1),\ldots,(D_m,\alpha_m)$ of matrices and strict angles.
Its finite-circle vector and vacuum coefficient are
\[
 h_N=A_N^{D_m}(\alpha_m)\cdots A_N^{D_1}(\alpha_1)\Omega_N,
 \qquad \langle\Omega_N,h_N\rangle.
\]
For two such words $h,g$, the proposed Gram form is
\[
 \langle h,g\rangle=\lim_{N\to\infty}\langle h_N,g_N\rangle.
\]
The following lemma establishes these limits. The empty word becomes
$\Omega$, and prefixing a row to the operator product defines its left
action. A bare vacuum coefficient means the same limit with the empty
word in the first slot; the adjoint identity expresses every Gram
coefficient as such a coefficient of a longer word.

\begin{lemma}[Plane kernels]\label{bulk:plane-kernels}
Finite words of homogeneous normalized strict-angle rows have
limiting Gram matrices as $N\to\infty$. Their quotient and completion
give a separable Hilbert space $\mathcal H$ with unit vector $\Omega$
and bounded left row actions satisfying
\begin{equation}\label{bulk:row-bound}
 \|A^D(\alpha)\|\le2\|D\|.
\end{equation}
Bare vacuum coefficients are continuous in strict angles and are
unchanged by a common translation of all time angles, as long as
the translated angles remain in $(-v,v)$.
\end{lemma}

\begin{proof}
We first identify the plane state, including the effect of the
zero-spin projection. Take periodic lists of time and space angles
in a strict range, with the two direction ranges ordered and
transversely separated in a half-circle. Use tori whose period
sizes are even multiples of the list periods. The resulting
rhombic square array is embedded and has bounded angles. Indeed
the cyclic order at each vertex gives local injectivity, while
averages of runs of either family have uniformly positive length
and uniformly transverse directions. Thus index distances control
Euclidean distances in both directions, giving properness and
global embedding.

Expand each $R$ into its two adjacent-side pairings with oriented
turn weight $e^{i\rho W/2}$. The pairing at the corner between the
two positive steps, and at its opposite corner, has coefficient
$a(u)$; the other pairing has coefficient $b(u)$. For example, the
mixed entry follows from
$a(u)e^{iu}+b(u)e^{-i(\lambda-u)}=1$ and its conjugate. In a
single-pairing entry the two turns cancel. A contractible loop
therefore has weight $d$. An essential simple oriented torus loop
has total turn zero. All essential loops are parallel up to
orientation. Their number $k$ is even, by the even intersection
parities with both period cycles and primitivity of their common
homology class. The allowed orientation count after projection
onto zero spin on either chosen quantum section is
\begin{equation}\label{bulk:essential-count}
 \chi=2^k\quad\text{if the intersection is zero},\qquad
 \chi=\binom{k}{k/2}\quad\text{otherwise}.
\end{equation}
Since $d<2$, there is $c(d)>0$ such that
$c(d)d^k\le\chi\le2^k$ for every even $k$.

Replacing $\chi$ by $d^k$ gives the canonical isoradial FK law,
up to the bounded torus genus factor of Lemma
\ref{prelim:loop-law}. To check the parameter, the relative
open weight at a primal corner between the positive steps is
$\sqrt q\,b/a$, and at the other primal corner it is
$\sqrt q\,a/b$. If $\gamma$ is the angle at the primal corner, the
odds divided by $\sqrt q$ are
\[
 \frac{\sin(\rho\gamma)}{\sin(\rho(\pi-\gamma))}.
\]
This is the canonical weight, with the complementary-angle
convention of \cite[Equation~(1.1), Theorem~1.1, and Equation~(1.3)]{DCLM2019Isoradial}. That theorem applies
to these fixed doubly periodic bounded-angle isoradial grids. It
gives uniqueness and buffered crossing estimates, uniformly over
the boundary partition on the outer buffer. The bulk passage-count
argument of Lemma \ref{prelim:passages} consequently applies
on each such grid. No estimate as an angle degenerates is used.

Here is why the essential-orientation tilt does not change the
local limit, even for unbounded aspect ratios. Let $r\le R$ denote
the shorter and longer period sizes. Cover the torus by $O(R/r)$
buffered boxes of size a small fraction of $r$, with a bounded
number of disjoint-buffer colors, so that every essential loop
crosses one of them. Include $O(R/r)$ translates, by medium periods,
of a target annulus about the query. Let $Y$ be the sum of the
passage counts in this cover. For each fixed $c>0$, conditional
passage moments, iteration within each color, and H\"older's
inequality give
\begin{equation}\label{bulk:torus-moment}
 \mathbb E e^{cY}\le e^{C_cR/r},\qquad k\le Y.
\end{equation}
The density of the tilted law is bounded above by $C e^{c_0Y}$
before normalization and its normalizing denominator is bounded
below. Consequently
$\mathbb E_{\rm tilt}e^Y\le C e^{CR/r}$, and Jensen's inequality
gives $\mathbb E_{\rm tilt}Y\le CR/r$. Translation invariance then
bounds the expected count at the target annulus uniformly.

Condition on the exterior of its inner box. The number of exterior
excursions from that boundary to distances comparable to $r$ before
returning is tight. If it is at most $K$, changing the filling
can alter only $O(K)$ essential loops: every affected loop must
leave the injective chart along such an excursion. The common
homology class can change only when no essential loop remains
unaffected. Thus the ratios of the genus and orientation tilt over
all fillings are bounded in terms of $K,q$. For the binomial case
this also follows directly by comparing central binomial
coefficients with indices differing by $O(K)$; a switch between
the two cases in \eqref{bulk:essential-count} can occur only when
there are $O(K)$ essential loops. The conditional inner law is
therefore comparable with FK under the induced partition.
Apply this density comparison once to the whole event that an arm or
medial strand from a fixed window reaches distance comparable to $r$.
On the event of at most $K$ exterior excursions, its conditional
probability is at most $C(K,q)$ times the corresponding FK bound.
Successive buffered circuits make that FK bound tend to zero as
$r\to\infty$. The excursion count is tight, so the complementary
probability becomes arbitrarily small by then letting $K\to\infty$.
No density factor is accumulated separately at each circuit scale.

Now condition outside a large fixed test box containing the query,
and let $z_0$ be its center. Choose a fixed $\eta>0$ so small that the
torus metric ball $B_T(z_0,2\eta r)$ has an injective planar lift;
the transverse period bounds above give this uniform choice. For large
$r$ the fixed box lies in $B_T(z_0,\eta r/2)$. The preceding escape
estimate shows that, with probability tending to one, no exterior
actual component or medial arc incident to the box reaches the rim of
$B_T(z_0,\eta r)$. Every such object is connected and meets the box,
so their entire union, together with every possible filling of the
box, lies in the same injective disk $B_T(z_0,2\eta r)$. Only exterior
components and arcs incident to the box can be changed by a new filling.
Every changed primal component is therefore planar, so its genus
contribution is unchanged; every changed medial loop is contractible.
The essential loop count, homology class, and orientation factor
therefore remain fixed. The conditional law of the filling is ordinary
FK. Letting the test box grow, comparison of increasing
cylinder events with free and wired laws and uniqueness identifies
the local plane law.

The same conclusion holds with finitely many marked matrices
$D_i$. Given the unoriented pairing, the absolute marked
orientation ratio is at most $\prod_i2\|D_i\|$: only marked loops
are broken, each newly cut arc has at most two orientations,
unaffected contractible loops retain weight $d$, and any residual
flux constraint has no more orientations than the central count
in \eqref{bulk:essential-count}. When all marked strands stay in
a fixed window, they do not encounter the distant flux projection.
The absence of escaping strands therefore proves marked local
convergence too.

Sending the time length to infinity first gives the
Perron-normalized coefficient at fixed spatial circumference.
Sending the circumference to infinity, and diagonalizing the two
limits, gives the stated plane coefficient. Any finite word can
be included in a periodic list; unmarked filler rows disappear
from its normalized longitudinal vacuum coefficient.

There is also a transverse quantization. Take a circle in the time
index, transfer in decreasing space index, and let the auxiliary
space spin flow in increasing time. The same difference relation
makes the untwisted transverse transfers commute and gives common
Perron data. The seam marks of the original word become one
tensor-product quantum insertion. If the time circle has length $L$,
write this insertion as $\mathcal D_L=\bigotimes_{t=0}^{L-1}D_t$,
where $D_t=I$ at unmarked times, using the seam-index convention of
\eqref{bulk:finite-row}. Let $\ell_L,r_L$ be the common left and right
zero-spin Perron vectors of the untwisted transverse transfers.
Sending the transverse cylinder length to infinity gives the normalized
coefficient
\[
 \frac{\langle\ell_L,\mathcal D_L r_L\rangle}
      {\langle\ell_L,r_L\rangle}.
\]
The transverse Perron eigenvalues cancel between the marked and
unmarked partitions. Their common eigenvectors, and hence this
coefficient, do not depend on the strict transverse angles. Sending
$L\to\infty$ gives the same plane coefficient as the longitudinal
limit already identified; this is why changing transverse angles
does not change it. Rotating both families and
then returning the transverse angles to $v$ proves common
time-angle translation invariance. Continuity follows from the
finite-box FKG sandwich and uniqueness as a common finite periodic
medium varies; sufficiently large fixed circuits at the limiting
medium control long marked strands, and all remaining local
orientation factors are continuous.

Finite-volume positivity now gives the limiting Gram form. To
prove boundedness before completing, let $h$ be a formal word
vector and set $T=(A^D)^*A^D$, with the adjoint understood first
as the reflected formal row in \eqref{bulk:adjoint}. The
marked expansion bounds
$\|T^jh\|^2\le C_h(2\|D\|)^{4j}$. Iterated Cauchy--Schwarz applied
to the positive form $\langle h,Th\rangle$ gives
\[
 \|A^Dh\|^2
 \le \|h\|^{2-2^{1-m}}
       \|T^{2^{m-1}}h\|^{2^{1-m}}.
\]
Letting $m\to\infty$ gives \eqref{bulk:row-bound}; the same inequality
shows that null vectors remain null. Continuity of kernels, and
words with rational angles and matrix entries, give separability.
\end{proof}

We may initially enlarge this construction by finite-position
quantum spin insertions and shifts. Moving a shift through a word
only translates finitely many markings, so the same proof gives
the enlarged Gram space and bounds. In that space the endpoint
rows are already defined by
\begin{equation}\label{bulk:endpoint-rows}
 A^D(v)=SD_0,\qquad
 A^D(-v)=(JD^tJ)_0S^{-1},
\end{equation}
where $(S\sigma)_j=\sigma_{j+1}$. These formulas follow from
$R(0)$ being the swap and \eqref{bulk:crossing}; their vacuum
eigenvalue is one. We next show that these endpoints preserve the
original cyclic space $\mathcal H$.

\subsection{A self-adjoint angle generator}

For a formal word $h$, translate all its angles by $t$ and call the
result $V_th$, for $t$ in a neighborhood of zero where the word is
strict. Reflection in its Gram formula negates the angles in the
reflected word. Common translation invariance therefore gives
\begin{equation}\label{bulk:hankel}
 \langle V_th,V_sg\rangle=\langle h,V_{s+t}g\rangle
\end{equation}
whenever both sides are defined. These curves are continuous. If
$h$ is null, then
$\|V_th\|^2=\langle h,V_{2t}h\rangle=0$ for small $t$, so the
translation is locally well defined on actual vectors.

\begin{lemma}[Angle generator]\label{bulk:boost}
There is a self-adjoint operator $B$ on $\mathcal H$ such that
$V_th=e^{tB}h$ locally for each word $h$. Ordinary strict-angle words
form a graph core for $B$. Each family $A^D$ extends to a bounded
holomorphic operator function on $|\operatorname{Re}z|<v$, with
\begin{equation}\label{bulk:imaginary-covariance}
 A^D(x+ib)=U_bA^D(x)U_{-b},\qquad U_b=e^{ibB}.
\end{equation}
\end{lemma}

\begin{proof}
For each word $h$, the function
$f_h(t)=\langle h,V_th\rangle$ is continuous and exponentially
convex on a sufficiently small interval: every permitted Hankel
matrix is a Gram matrix by \eqref{bulk:hankel}. The interval
Bernstein--Widder theorem, in the form of
\cite[Theorem~3.4 and Corollary~3.5, p.~735]{BuescuPaixao2018}, represents
$f_h$ as a Laplace transform of a positive measure with finite
integrals throughout that interval. In particular it is analytic
near zero. Inner products of finite differences of $V_th$ are
the corresponding differences of $f_h(s+t)$. Their limits show
that all Hilbert-space derivatives $h^{(k)}$ exist and satisfy
\begin{equation}\label{bulk:analytic-vector-bound}
 \|h^{(k)}\|^2=f_h^{(2k)}(0)
 \le C_h(2k)!r_h^{-2k}.
\end{equation}
The Taylor series recovers the vector curve: this follows by
comparing its inner products and its squared norm using
\eqref{bulk:hankel}.

On the span of all these derivative vectors, define
$B_0h^{(k)}=h^{(k+1)}$. The cross versions of
\eqref{bulk:hankel} make this well defined and symmetric. More
explicitly, if a linear combination of derivatives vanishes, its
proposed derivative has zero inner product with every word by
differentiating the cross identity; density makes that derivative
zero. The same identity gives symmetry between derivative
vectors. The domain is dense and invariant. Bound
\eqref{bulk:analytic-vector-bound}, also with $k$ replaced by $k+j$,
implies that every vector in this domain is analytic for this same
operator: for some $t>0$ its series
$\sum_{j\ge0}t^j\|B_0^jh^{(k)}\|/j!$ converges. Nelson's analytic-vector theorem \citep[Lemma~5.1]{Nelson1959}, in the
form stated in \citet[Definition~4.7 and Theorem~4.8]{Neeb2011}, makes $B_0$
essentially self-adjoint. Denote its closure by $B$.

Finite differences of ordinary words approximate $h^{(k)}$ both
in norm and after applying $B$, by the Taylor series just obtained.
Ordinary words are thus a graph core. The vector series gives
$V_th=e^{tB}h$ for small real $t$ on each word.

For fixed $D$, the row family is strongly continuous, by the Gram
continuity and uniform bound. The word identity
$A^D(x+t)V_th=V_t(A^D(x)h)$ yields, weakly on the graph core,
\[
 \frac{d}{dx}\langle g,A^D(x)h\rangle
 =\langle Bg,A^D(x)h\rangle-\langle g,A^D(x)Bh\rangle.
\]
Integrate this equality and approximate $g,h$ in graph norm to
extend it to $\operatorname{Dom}B$. The extension in
Equation~\eqref{bulk:imaginary-covariance} then satisfies the Cauchy--Riemann
equations between domain vectors, with continuous partial
derivatives. Its uniform bound $2\|D\|$ extends weak holomorphicity
from that dense set to all vectors, hence gives bounded operator
holomorphicity throughout the strip.
\end{proof}

Words whose angles all lie in the middle third $(-v/3,v/3)$ already
span densely: test orthogonality to them and apply analyticity
successively in every argument of a word. For a fixed twist,
analytic continuation of the commuting-row identity and
\eqref{bulk:adjoint} gives
\begin{equation}\label{bulk:P-star}
 [P_\zeta(z),P_\zeta(w)]=0,\qquad
 P_\zeta(z)^*=P_\zeta(-\overline z).
\end{equation}
Thus this family lies in an abelian von Neumann algebra.  Decompose
its representation into cyclic summands, or equivalently into a direct
integral with multiplicity spaces: each algebra element acts on a fiber
by a scalar times the identity, while Hilbert-space vectors remain
vector-valued on that fiber.  The cyclic representation and spectral
calculus are recalled in \citet[Sections~3.8 and 4.3]{Peterson2020OperatorAlgebras}.
Its simultaneous spectral variables can therefore be chosen as scalar functions
$f(z)$, analytic almost everywhere on the strip, bounded by $2$,
and satisfying $f(-\overline z)=\overline{f(z)}$. To obtain one
common exceptional set, start with a multiplication representation
of the algebra generated by countably many arguments in a dense
set, and use the norm-bounded Taylor series on countably many
discs. Covariance \eqref{bulk:imaginary-covariance} says that
$f(z)\mapsto f(z+ib)$ preserves the spectral measure class.

\subsection{Continuation through the two edges}

To estimate long products of rows, we next need uniform control of
their scalar spectral functions up to the two boundary lines of the strip.

\begin{lemma}[Scalar continuation in a fixed band]\label{bulk:edge-band}
Almost every spectral function $f$ extends analytically and
without zeros to bands of one fixed positive width about the
lines $\operatorname{Re}z=\pm v$. It has modulus one on those lines
and satisfies $|f(z)|\le1$ in the intervening strip. At any specified
point of the strip, the set of spectral functions vanishing there
has measure zero.
\end{lemma}

\begin{proof}
We first obtain a uniform finite-volume vacuum estimate. Temporarily
write $T_D(u)=T_N^D(v-u/\rho)$. For $u$ in a fixed sufficiently
small complex disc about zero, both
\begin{equation}\label{bulk:division}
 T_D(u)T_I(u)^{-1},\qquad T_I(u)^{-1}T_D(u)
\end{equation}
are holomorphic with norms at most $C\|D\|$, uniformly for all
sufficiently large $N$.

To prove the first assertion, multiply by the transfer
$\widetilde T_I(-u)$ having reversed auxiliary order. Fold its
auxiliary index $c'$ by transposition. The product is the
two-auxiliary partial trace of
$(D\otimes I)\prod_jG_j$, in increasing site order, where
\[
 G_j=R_{cj}(u)R_{c'j}(-u)^{t_{c'}}.
\]
The line spanned by $|++\rangle+|--\rangle$ is invariant in the
auxiliary space, and on this line $G_j$ is
$r(u)I$, where $r(u)=1-b(u)^2$, by inversion. At $u=0$ the range
of $G_j$ is entirely in that line tensored with the site, by the
swap identity. In the common triangular decomposition, $G_j/r(u)$
has identity on this line, a complementary diagonal block of
norm at most $C|u|$, and a bounded off-diagonal block. A product of
these matrices has uniformly bounded norm, since the
off-diagonal block is bounded by a geometric series. Its partial
trace without $D$ differs from the identity by at most
$C(C|u|)^N$: the off-diagonal block has zero partial trace, the
line block gives the identity, and the remaining diagonal product
has the stated norm. This partial trace is invertible on one
fixed disc, proving the first bound in \eqref{bulk:division}.
Reversing the product in $G_j$ gives the same argument with the
invariant covector, proving the second. Crossing gives the
corresponding bounds near $u=\lambda$.

The common vacuum is an eigenvector of $T_I(u)$ by analyticity,
and its eigenvalue is nonzero in these discs. Hence
$A_N^D(z)\Omega_N$ has a uniform vector bound in discs of a fixed
radius about $z=\pm v$. Pairing with each fixed limiting word,
normal-family compactness gives analytic subsequential limits;
their values on the interior fix them uniquely. The vector norm
bound identifies these limits with a bounded holomorphic
$\mathcal H$-valued function. To identify its endpoint values, Cauchy's
derivative estimate on a smaller fixed disc gives, uniformly in $N$,
\[
 \bigl\|(A_N^D(x)-A_N^D(e))\Omega_N\bigr\|
 \le C_D|x-e|,\qquad e\in\{-v,v\},
\]
for real interior $x$ sufficiently close to $e$. The enlarged Gram
construction includes the endpoint swap vectors from
\eqref{bulk:endpoint-rows} and their mixed scalar products with strict
row words. Passing to those Gram limits preserves the squared norm
inequality. Letting $x\to e$ therefore identifies the continued vacuum
vector with the swap vector strongly in the enlarged space. This is
convergence in the angle, not a claim of strong convergence between
different finite-circle spaces.

Now commute this endpoint row through a word with all angles in
$(-v/3,v/3)$. Each exchange uses an intertwiner whose inverse is
analytic on one fixed disc about the endpoint, independent of the
word: the middle-third angle restriction keeps its parameter
uniformly away from inversion zeros. The exchange expresses the
result as a finite sum of bounded interior rows applied to the
continued vacuum vectors. It follows that the endpoint preserves
$\mathcal H$ and that $A^D(z)h$ continues holomorphically on this
same disc for every middle-third word $h$. Its bound may depend
on the word, because the number of exchanges does. This does not
alter the radius. Strong convergence of interior rows to the
endpoint holds first on this dense set and then on all of
$\mathcal H$, by the uniform interior operator bound.

Choose a countable dense set of middle-third words $h_j$, and a
finite equivalent measure in the spectral representation of
$P_\zeta$. If $r$ is the common continuation radius, write the
continued vector as $F_j(z)=\sum_{k\ge0}a_{jk}(z-v)^k$ at the
right endpoint. For $r'<R<r$, Cauchy's estimate gives
\[
 \|a_{jk}\|_{L^2}\le M_{j,R}R^{-k},\qquad
 \sum_{k\ge0}(r')^k\|a_{jk}\|_{L^2}<\infty.
\]
Since the measure is finite, the corresponding pointwise series
of fiber norms is summable almost everywhere. Intersect these
full-measure sets over the countable core and a countable family
of smaller radii. We obtain pointwise holomorphic vector series
on the same disc for every $j$. The left endpoint is identical.
On its interior portion the vector equals $f(z)h_j$.
Analyticity consequently makes the continued fiber proportional
to $h_j$ everywhere. At almost every spectral point at least one
vector in the countable dense set is nonzero; ratios using any
such vector agree on the interior and hence throughout the disc.
This constructs the scalar continuation with one common radius.

The swap formulas show that $P_\zeta(v)$ and $P_\zeta(-v)$ are
unitary inverses. Strong boundary convergence and
\eqref{bulk:imaginary-covariance} give unitary horizontal limits
at $\pm v+ib$. For countably many $b$ within the continuation
discs, the scalar value there has modulus one almost everywhere.
Continuity gives modulus one on each full vertical diameter.
Schwarz reflection across that diameter now implies
$f(z)\overline{f(2v-\overline z)}=1$, and the analogous identity
at $-v$, throughout the symmetric disc. Thus there are no zeros
there. Rational imaginary translations preserve null sets and
cover both boundary lines by discs of the same radius. This gives
the asserted fixed bands.

The bounded-function maximum principle on the strip gives
$|f|\le1$. Finally each nonzero analytic $f$ has isolated zeros,
so Fubini's theorem shows that on a fixed vertical orbit almost
every argument has no zero for almost every spectral function.
The covariance of the spectral measure class transfers this
null-set assertion to any prescribed point of that orbit.
\end{proof}

\subsection{Inner factorization and low spectral bands}

Use $w=e^{iz}$ to map $|\operatorname{Re}z|<v$ to the right
half-plane. We keep the notation $f$ for the resulting function.
Lemma \ref{bulk:edge-band} says that $f$ is inner, continues through
every finite nonzero point of the imaginary axis, and that all its
zeros satisfy
\begin{equation}\label{bulk:zero-cone}
 \operatorname{Re}w_j\ge c|w_j|
\end{equation}
for one $c>0$ independent of the spectral point. The disk
factorization and its boundary-spectrum description
\cite[pp.~1367--1368, Equations (1)--(2)]{AiDongJiang2019}, transported
by $(w-1)/(w+1)$, therefore give
\begin{equation}\label{bulk:inner-factorization}
 f(w)=\eta e^{-aw-b/w}\prod_j\beta_j(w),\qquad
 a,b\ge0,\quad |\eta|=1.
\end{equation}
Here $\beta_j$ is $(w-w_j)/(w+\overline{w_j})$ with the canonical
constant phase, chosen positive at $w=1$ when $w_j\ne1$.
The zeros satisfy the half-plane Blaschke condition
$\sum_j\operatorname{Re}w_j/(1+|w_j|^2)<\infty$.
There is no other singular factor: analytic continuation across
the rest of the boundary excludes singular measure there, leaving
only atoms at $0$ and $\infty$. Under the Cayley map these two
atoms give exactly $b/w$ and $aw$.

Suppose $-\log|f(1)|\le s$, where $s>0$ is sufficiently small.
There is then no zero at $1$, and evaluation of
\eqref{bulk:inner-factorization} gives
\begin{equation}\label{bulk:zero-budget}
 a+b+\frac12\sum_j
 \log\left(1+\frac{4\operatorname{Re}w_j}{|1-w_j|^2}\right)\le s.
\end{equation}
The cone condition implies, with a fixed constant,
\begin{equation}\label{bulk:small-large}
 a+b+\sum_j\min\{|w_j|,|w_j|^{-1}\}\le Cs.
\end{equation}
Indeed, for $r=|w_j|$, the fraction inside the logarithm is at least
$4cr/(1+r)^2$, and the resulting logarithm is bounded below by a
positive constant times $\min(r,r^{-1})$.

Choose $c_1>0$ sufficiently small. On the half-annulus
\begin{equation}\label{bulk:annulus}
 c_1^{-1}s<|w|<c_1/s,\qquad |\arg w|\le\pi/2,
\end{equation}
$f$ and $f^{-1}$ are bounded by fixed constants. On the unit
semicircle,
\begin{equation}\label{bulk:near-sign}
 f(w)=f(1)+O(s).
\end{equation}
To see both assertions directly, \eqref{bulk:small-large} separates
every zero into $|w_j|\le Cs$ or $|w_j|^{-1}\le Cs$. For a small
zero the ratio $\beta_j(w)/\beta_j(1)$ differs from one by at most
$C(1+|w|^{-1})|w_j|$; for a large zero its error is at most
$C(1+|w|)/|w_j|$. Their logarithms are summable and uniformly
bounded on \eqref{bulk:annulus}, after decreasing $c_1$.
The exponential factor is controlled by
$a|w-1|+b|w^{-1}-1|$. On $|w|=1$ the same estimates are $O(s)$.
This proves the claims, including bounded invertibility.

For a twist $\zeta$, define the nonnegative spectral operator
\[
 H_\zeta=-\log|P_\zeta(0)|,\qquad
 E_\zeta(s)=\mathbf1_{[0,s]}(H_\zeta).
\]
The no-zero assertion of Lemma \ref{bulk:edge-band} makes this
definition unambiguous almost everywhere. Since $P_\zeta(0)$ is
self-adjoint, split $E_\zeta(s)$ further according to its sign,
denoting these projections by $E_\zeta^\sigma(s)$,
$\sigma=\pm1$. Equations \eqref{bulk:annulus}--\eqref{bulk:near-sign}
say that the compressed $P_\zeta$ is boundedly invertible on the
half-annulus and
\begin{equation}\label{bulk:sign-approx}
 P_\zeta(w)E_\zeta^\sigma(s)
 =\sigma E_\zeta^\sigma(s)+O(s)
 \quad (|w|=1,\ \operatorname{Re}w\ge0).
\end{equation}

\subsection{Reversal, Laurent modes, and the sandwich}

The low-band estimates now control the surrounding twisted rows.  It
remains to identify the leading angular mode of $N$ and the improved
decay obtained by canceling it in \eqref{bulk:high-estimate}.

The endpoint formulas give the exact reversal identity
\begin{equation}\label{bulk:reversal}
 M(v)+P_\xi(v)M(-v)P_\zeta(v)=0.
\end{equation}
For example, if $D=L_-$, then
$K(\xi)DK(\zeta)=(\zeta/\xi)D$ and
\[
 e^{-i\tau\rho v}(\zeta/\xi)=-e^{i\tau\rho v}.
\]
The raising case is the analogous computation, with
$\tau=-\epsilon$. The identity at $-v$ follows using the inverse
endpoint twists. Imaginary covariance extends these relations
after every common imaginary shift: both $M$ terms acquire the
same scalar multiplier. Boundary convergence is strong.

Compress $M$ between the low bands of the two twists:
\[
 M_c(w)=E_\xi(s)M(w)E_\zeta(s).
\]
For $\operatorname{Re}w<0$ define it from the right half-plane by
\begin{equation}\label{bulk:gluing}
 M_c(-w)=-P_\xi(w)^{-1}M_c(w)P_\zeta(w)^{-1},
 \qquad \operatorname{Re}w>0,
\end{equation}
where the inverses are only on the stated low bands. The two
versions of \eqref{bulk:reversal} give equality on both imaginary
rays. Uniform bounds and strong boundary continuity allow weak
holomorphic gluing, for instance by Morera's theorem on small
rectangles crossing the rays. Thus $M_c$ is holomorphic on the
full annulus \eqref{bulk:annulus}. Throughout it,
\begin{equation}\label{bulk:Mc-bound}
 \|M_c(w)\|\le C|w|^{\tau\rho}.
\end{equation}

Write its Laurent expansion as $\sum_{j\in\mathbb Z}w^jM_{c,j}$.
Cauchy estimates on circles near the inner or outer edge give
\begin{equation}\label{bulk:Laurent}
 \|M_{c,j}\|\le C(C's)^{|j-\tau\rho|}.
\end{equation}
On $|w|=1$, summing these estimates yields
\begin{align}
 \|M_c(w)\|&=O(s^\rho),\label{bulk:Mc-small}\\
 \frac{M_c(w)-M_c(-w)}2
 &=w^\tau M_{c,\tau}+O(s^{1+\rho}),\qquad
 \|M_{c,\tau}\|=O(s^{1-\rho}).\label{bulk:odd-mode}
\end{align}
Here the nearest odd integer to $\tau\rho$ is $\tau$; the next
possible distance is $1+\rho$. This is where $0<\rho\le1/3$ is
used in the ordering of the relevant exponents.

Refine the left and right low bands by signs $\sigma_\xi,\sigma_\zeta$.
If the signs agree, \eqref{bulk:gluing} and
\eqref{bulk:sign-approx} imply
$M_c(-w)=-M_c(w)+O(s)\|M_c(w)\|$ on the unit semicircle.
By \eqref{bulk:Mc-small}, its even part is $O(s^{1+\rho})$.
On this block,
\[
 E_\xi^{\sigma_\xi}(s)N(w)E_\zeta^{\sigma_\zeta}(s)
 =2\sigma_\xi w^\tau M_{c,\tau}+O(s^{1+\rho}).
\]
If the signs disagree, the two leading signs in
$M_cP_\zeta+P_\xi M_c$ cancel, giving $O(s^{1+\rho})$ directly.
Consequently, on the full low bands,
\begin{align}
 \|E_\xi(s)N(x)E_\zeta(s)\|&\le Cs^{1-\rho},
 \label{bulk:low-band}\\
 \|E_\xi(s)\{N(x)-e^{i\tau(x-y)}N(y)\}E_\zeta(s)\|
 &\le Cs^{1+\rho}.\label{bulk:high-band}
\end{align}
All real $x,y\in[-v,v]$ correspond to the closed unit semicircle.
A product of an even bounded number of extra rows is $I+O(s)$ on
each sign block, by \eqref{bulk:sign-approx}. Its error multiplied
by \eqref{bulk:Mc-small} is $O(s^{1+\rho})$. This proves exactly
the same conclusions with the extra pairs stated in Theorem
\ref{bulk:angular-estimate}, even if the two terms use different
pairs.

To finish the sandwich estimate, decompose each $H_\zeta,H_\xi$
into $[0,n^{-1}]$, dyadic intervals above $n^{-1}$ up to a small
fixed threshold, and the remaining interval. Schwarz--Pick gives,
uniformly for $x$ in the fixed compact interval $I$,
\[
 1-|f(x)|\ge c_I(1-|f(0)|).
\]
One can obtain this directly from the pseudohyperbolic contraction
between $0$ and $x$; their hyperbolic distance stays bounded on
$I$. A band with lower endpoint $u$ below the fixed threshold is
therefore damped by the $n$ strict steps by $e^{-cnu}$; the
remainder is damped exponentially in $n$. All other rows are
contractions, including the endpoint rows. Apply
\eqref{bulk:low-band} or \eqref{bulk:high-band} with $s$ the larger
of the two band upper endpoints. The double dyadic sum of
\[
 C\max(u,u',n^{-1})^\beta e^{-cn(u+u')},
 \qquad \beta=1-\rho\ \text{or}\ 1+\rho,
\]
is at most $C_\beta n^{-\beta}$. Terms involving a remainder band
use the uniform unrestricted row bound and are exponentially
small. This proves \eqref{bulk:low-estimate} and
\eqref{bulk:high-estimate} in the homogeneous representation.

\subsection{Diagonal cuts and the physical reflection norm}

It remains to prove the final assertion of Theorem
\ref{bulk:angular-estimate}. Work in the diagonal time frame and
write
\[
 \gamma_0=\pi/4,\quad X=e^{-i\gamma_0},\quad
 Y=e^{i\gamma_0},\quad H=-X,\quad L=Y.
\]
Space steps alternate $H,L$, starting with $H$; time steps are
$X$ or $Y$. Put $R_*=R(\lambda/2)$ and let
\begin{equation}\label{bulk:G}
 G=\prod_j\bigl(\operatorname{swap}R_*\bigr)_{2j,2j+1}.
\end{equation}
The factors act on disjoint pairs. Each is positive definite:
its middle block has eigenvalues $1\pm b(\lambda/2)>0$, and the
two outer entries are $a(\lambda/2)>0$. Use $G$ as the
finite-circle metric. Directly from the collapsed tiles,
\begin{equation}\label{bulk:diagonal-rows}
 T_N^D(\gamma_0)=D_{N-1}SG,\qquad
 T_N^D(-\gamma_0)=(JD^tJ)_0S^{-1}G.
\end{equation}

Here is the geometry behind both the metric and these formulas.
The lower slice has vertices $p,p-X,p+Y-X,\ldots$.
Reflection in the cut through $p$, perpendicular to time, gives
upper order $L,H$. For lower spins $(c_0,c_1)$ and upper spins
$(d_0,d_1)$, the intervening square kernel is
\[
 (R_*)_{-c_0,d_0;-d_1,c_1}
 =(\operatorname{swap}R_*)_{d_0,d_1;c_0,c_1}.
\]
Shifting the slice origin by $Y$ or $X$ puts the seam at site $-1$
or $0$ respectively, with reversed conventions in the latter
case. Equivalently, every second indexed tile has collapsed.
A coincident-step tile gives $R(0)$, the swap; an opposite-step
tile gives $R(\lambda)$, which imposes
$h_0=-s_1,\ h_1=-s_0$, with weight one. Its duplicate sides are
therefore spin identifications. The adjoint for the $G$ metric
exchanges the two time steps and sends $D$ to $J\overline D J$.

Both untwisted operators $S^{\pm1}G$ are primitive on the
zero-spin sector. Over a complete shift cycle, paired bonds of
both parities are available. One can hold every pair or swap a
chosen adjacent opposite pair, and such swaps connect the sector.
The transfers commute, have a common positive right Perron
vector, and have the same eigenvalue: their relation by $S^2$
and uniqueness make the vector $S^2$-invariant. Their left
Perron data are its $G$-dual. Normalize the vacuum to unit
$G$ norm.

These degenerate indexed arrays are exactly physical square
tori; they need no limiting estimates for degenerate rhombi.
To check this in both quantizations, let $p_t$ follow a periodic
sequence of $X,Y$ steps with even period size, and set
\[
 q_{2j}=j(Y-X),\qquad q_{2j+1}=q_{2j}-X.
\]
The uncollapsed tiles of the array $p_t+q_j$ biject with the
physical squares centered at each time level. The collapsed
identifications join duplicate physical edges as just described.
There are no free cycles consisting solely of duplicates:
each collapsed pair uses a time edge, which by the alternating
space pattern is also incident to an uncollapsed tile. Hence
zero-spin projection on either an unmarked zigzag slice or a
periodic time path is precisely the corresponding physical flux
projection. Marks at the time seam mark that physical path.
The periods preserve parity and are uniformly transverse, so
the torus-to-plane argument of Lemma \ref{bulk:plane-kernels}
applies directly to these square kernels.

For clarity, the transverse period product is primitive too.
Prescribe arbitrary zero-sum spin sequences on the time paths
$p_t$ and $p_t+m(Y-X)$. Interpret them as increments of integer
heights with consistent sublattice parity, periodically bounded
on the two lifted paths. Each path is geodesic in the square
graph with steps $X,Y$, so its own heights are $1$-Lipschitz.
Between the two paths the graph distance is at least $|m|$:
if the time-level difference is $\ell$, the two lower bounds
are $|\ell|$ and $2|m|-|\ell|$. Thus for all sufficiently large
$|m|$ the prescribed heights together are $1$-Lipschitz.
Extend them by the minimum of the integer distance cones
$h(a)+\operatorname{dist}(a,\cdot)$ from prescribed vertices $a$.
This extension is time-periodic; its parity is fixed, so every
edge increment is exactly $\pm1$. It gives an allowed spin
configuration of positive weight across the transfer strip.
Therefore a sufficiently large power of either oriented
two-space-step transfer has all entries positive. The
difference relation makes this period product commute with
the untwisted transfer at strict transverse angle $v$ on the
same time chain; the intertwiners involved are invertible.
They consequently share Perron data. Seam-marked transverse
vacuum coefficients agree with those at that strict angle.

Construct the limiting Gram space with the $G$ metric, allowing
the two endpoint rows, local site matrices, and $S^2$.
The adjoint of a site matrix, conjugated by $G$, is still local
and bounded. The preceding physical cylinder limits and the
marked moment argument give the Gram limits and bounded row
actions. Let $\mathcal H_*$ be the closed subspace generated
by seam words in $A^D(\pm\gamma_0)$ alone. Their Gram kernels
are the homogeneous transverse kernels, by the adjoint formula,
the transverse Perron comparison, and then the plane limit.
Thus $\mathcal H_*$ identifies isometrically with the subspace
of $\mathcal H$ generated by those words.

This subspace also contains the physical states required in
the theorem. It is reducing for $P_1(\pm\gamma_0)$. Those
operators have dense range there, because their adjoints have
no kernel by the no-zero assertion at the specified interior
angles. The finite-row identities give
\[
 S^2P_1(-\gamma_0)=P_1(\gamma_0).
\]
Both shifts by two therefore preserve $\mathcal H_*$, first
on a dense range and then by continuity. Equation
\eqref{bulk:diagonal-rows} with arbitrary $D$ similarly shows
that site matrices at $-1$ and $0$ preserve it; shifting by two
gives every finite site. Every physical edge between tiles
lies on a time-level zigzag slice, so arbitrary finite physical
modifications are such site insertions interleaved with layer
transfers. They belong to $\mathcal H_*$. For ordinary
homogeneous slicing the same assertion follows from the
endpoint shifts; several time-edge insertions on a strict
row can also be moved to one seam by commuting through finite
strings of invertible $R$ matrices and expanding their matrix
elements.

Finally, reflection in a time-level cut, followed by complex
conjugation and arrow reversal, matches lower spins to upper
spins through the intervening kernel $G$. Normalize by the
unmarked vacuum, or equivalently by the unit $G$ vacuum and
one common Perron factor per layer. The squared norm of a
lower state is then exactly the vacuum expectation ratio
of its reflected configuration. Between two cuts, transfer
from the lower slice of the first to the lower slice of
the second, and glue to the reflected future there using
$G$. Strings immediately below and above the last cut act
on the respective states and leave this kernel unchanged.
The homogeneous estimate restricted to $\mathcal H_*$ is
therefore the asserted diagonal estimate. This completes the
proof of Theorem \ref{bulk:angular-estimate}.

\section{Local path calculus for the current}\label{sec:current}

The transfer estimate applies to a lattice current because its
path dependence can be removed locally. This section proves the
needed tensor identities and records the phase bookkeeping when
a current is placed between two diagonal cuts.

\subsection{A path primitive and its deformation rule}

Direct a simple path $\mathcal C$ of tile edges from its initial
point. Measure each crossing spin using the left normal of the
path step. At a crossing, a matrix has row index on the right tile
and column index on the left tile. Lift the tangent angle
continuously along a rounded version of the path, starting with a
specified lift. Define its primitive $\mathcal J(\mathcal C)$ by
summing over the choice of one marked step:
\begin{itemize}
\item at that step insert $e^{i\rho\alpha}L_-$, where $\alpha$
      is the lifted tangent angle;
\item on all preceding steps insert $K(Q)$;
\item on all later steps insert the identity.
\end{itemize}
Thus both strand arrows at the marked crossing point away from it.
All insertions are additional to the tile tensors. Other primitives
and other marked insertions are kept outside the deformation
corridor, unless their order is explicitly included.

\begin{lemma}[Local deformation]\label{cur:deformation}
In the positive winding expansion, the primitive is invariant,
pairing by pairing, under simple-path isotopies that continue its
initial lift, keep its endpoints in their regions, and avoid the
other insertions. The statement is a tensor identity for arbitrary
outside spins. A pure diagonal string obeys the corresponding
deformation rule without a source.
\end{lemma}

\begin{proof}
Draw the strands smoothly and disjointly, with transverse path
crossings. At a marked crossing use the tangent-angle lift rotated
by less than $\pi/2$ to the orthogonal crossing direction; this
agrees with the lattice convention at a tile edge. Sliding the
source on its strand while its tangent changes by $b$ changes the
sum of the turns of the two departing pieces by $-2b$.
The winding factor is consequently multiplied by $e^{-i\rho b}$,
which cancels the change of the marker gauge. All other spin
constraints and prefix factors correspond directly.

A generic isotopy also creates or removes a pair of crossings at
a tangency. Terms marked elsewhere are unchanged because the
two new diagonal factors cancel. The two terms marked at the
new crossings cancel one another. To check the sign, orient a
test tangent $\beta$ of the strand approximately in the path
direction near tangency, so crossings $1,2$ occur in that order
both along the path and along this test direction. With the
source at crossing $2$, let $c_1$ be the signed spin at crossing
$1$; the strand there travels opposite to $\beta$, away from the
source. The orthogonalized lifts satisfy
\begin{equation}\label{cur:tangency-lift}
 (\alpha_2-\beta_2)-(\alpha_1-\beta_1)=-c_1\pi.
\end{equation}
For example, when the first path tangent is clockwise from
$\beta$, the offsets are $-\pi/2,+\pi/2$ and $c_1=-1$; the other
case reverses the signs. The ratio of the marker gauge times
the source-piece winding factors is therefore $e^{-ic_1\lambda}$.
The extra prefix factor is $Q^{c_1}$, and
\[
 e^{-ic_1\lambda}Q^{c_1}=-1.
\]
The outside constraints and factors of the two terms are
identical. This proves the cancellation.

A generic isotopy has only isolated tangencies of this kind,
besides transverse slides, and can be chosen to keep all fixed
insertions disjoint. This proves the primitive identity. For a
pure diagonal string, only the first cancellation is needed;
equivalently it is the tile ice-conservation identity.
\end{proof}

\subsection{Boundary-tree coefficients}

We use the boundary factors from Section \ref{sec:exponents}.
For reference, if $\theta_j$ is the lifted inward normal at port
$j$, $h_j$ its real height parameter, and $c=\pm1$ is its
inward/outward arrow sign, they are
\[
 b_j(c)=(2\sin\lambda)^{-1/2}
       \exp\{ic(\rho\theta_j/2+\lambda h_j)\}.
\]
Consider a disk, order its boundary from a flat gap $p$, and
assume there are no obstructing extra strings. Start a path from
$p$ toward an interior vertex $z$ with initial tangent lift
strictly between the positive boundary tangent and that tangent
plus $\pi$. Every boundary arc bounds a child interval not
containing $p$; these arcs form the separator tree of the
complementary regions.

\begin{lemma}[Tree evaluation]\label{cur:tree-weights}
For an arc $e:i\to j$ bounding its child interval, put
$m=h_i-h_j$ and
\begin{equation}\label{cur:weights}
 \begin{aligned}
 w_e&=\frac{\cos((m-\tfrac12)\lambda)}{\sin\lambda},&
 k_e&=-\frac{\cos((m+\tfrac12)\lambda)}{\sin\lambda},\\
 a_e&=e^{-i\lambda(h_i+h_j)},&
 c_*&=\frac{e^{-i\lambda}}{2\sin\lambda}.
 \end{aligned}
\end{equation}
For each fixed pairing, the primitive to $z$ is the sum over a
marked edge of the separator-tree path from $p$ to the region of
$z$: use $k_e$ on earlier edges, $c_*a_e$ on the marked edge, and
$w_e$ on every other boundary arc. Unmarked internal loops have
weight $d$, and a marked closed loop contributes zero.

For a pure $K(Q)$ string to $z$, use $k_e$ on the full
separator-tree path and $w_e$ off it. An internal loop enclosing
$z$ then has weight $-2$ instead of $d$. The parity of the number
of such loops records the color change from the boundary-tree
region to the lattice color of $z$.
\end{lemma}

\begin{proof}
By Lemma \ref{cur:deformation}, for each pairing we may choose a
simple path crossing minimally the boundary arcs separating $p$
from $z$, followed, if necessary, by the nested closed loops.
Such paths exist in the disk regions cut out by the disjoint
arcs and loops. A sole outward source on a closed loop is
inconsistent with the remaining arrow constraints, so it gives
zero.

For a parent-to-child crossing, the boundary-ordered arc travels
along the positive normal of the path. Let $\beta$ be its
tangent lift, continued from $\theta_i$. The correct lift relation
is the exact identity
\begin{equation}\label{cur:exact-angle}
 \beta-\alpha=\pi/2,
\end{equation}
without a multiple of $2\pi$. To verify this, close a simple curve
by taking the boundary from $p$ to $i$, the arc to the crossing,
and the reversed path back to $p$. Its positive full turn, with
the left right-angle turns at the port and crossing and the
short turn at $p$, is $\beta-\alpha+3\pi/2$. It equals $2\pi$,
proving \eqref{cur:exact-angle}.

For an outward source the sum of strand turns is
$\theta_i+\theta_j-\pi-2\beta$. Multiplying its winding weight by
the two outward boundary factors and by $e^{i\rho\alpha}$ gives
\[
 \frac1{2\sin\lambda}
 e^{-i\lambda(h_i+h_j)}
 e^{i\rho(\alpha-\beta-\pi/2)}
 =c_*a_e.
\]
Summing the two orientations of an unmarked arc gives $w_e$.
If a prefix string crosses it, multiplication by $K(Q)$ changes
that sum to $k_e$. These are precisely the tree weights in
\eqref{cur:weights}.

For a pure string, the same calculation holds on every separating
boundary arc. A counterclockwise enclosing loop has intersection
factor $Q^{-1}$, so the sum of its two orientations is
$e^{i\lambda}Q^{-1}+e^{-i\lambda}Q=-2$.
Crossing each such medial loop switches the adjacent color,
which proves the final parity assertion.
\end{proof}

\subsection{Reflection and two-step increments}

Reflect in a line of angle $\phi$, and combine this operation
with complex conjugation and arrow reversal. Relative to the
mirrored step direction, spins are unchanged, while the left and
right tiles are exchanged. Consequently the reflected primitive
has the form
\[
 L_+,\ K(Q^{-1}),\qquad e^{i\rho\alpha'},
 \qquad \alpha'=2\phi-\alpha,
\]
up to a common scalar. It obeys the reflected local deformation
identities, using the opposite winding expansion. In the notation
of \eqref{bulk:charges}, it has $\epsilon=-1$ and still $\tau=1$.
Further reflections alternate between these two types. Moreover,
a high difference with coefficient
$e^{i(\arg d_1-\arg d_2)}$ becomes the same expression with the
reflected directions $d'_1,d'_2$.

A \emph{low stencil} is a primitive increment on a straight
two-step extension from a vertex $z$. A \emph{high stencil} is
the difference of two such increments in distinct tile-axis
directions $d_1,d_2$, with coefficient
$e^{i(\arg d_1-\arg d_2)}$ on the subtracted term.
Both extensions have one common arriving prefix with fixed lift.
The local path choices are simple and extensible.

\begin{proposition}[Extraction in a clear slab]\label{cur:slab-estimate}
Suppose a diagonal-time slab of macroscopic thickness about a
low or high stencil contains no other non-diagonal tensors or
endpoints. Diagonal chains may cross the slab. The stencil has an
intact neighborhood disjoint from all other chains except its
own prefix, and the slab has spare margins for routing and cuts.
Bounded lattice errors are permitted inside these margins.
Then the stencil contraction can be written as the corresponding
low or high expression in Theorem \ref{bulk:angular-estimate},
between exterior states independent of the choice of stencil
term. Only boundedly many compensating pairs of rows and a
common bounded scalar are introduced.

In particular, if the two clear parts of the slab contain at
least $n$ physical steps, its scalar contraction is bounded by
$Cn^{-1+\rho}$ for a low stencil, and by $Cn^{-1-\rho}$ for a
high stencil, times the two exterior reflection norms. Reflected
stencils satisfy the same statement.
\end{proposition}

\begin{proof}
There are two matters to check: the compared paths must have
common endpoints, and rerouting the other charges must contribute
the same phase to every possible source position.

Let $a,b$ be the two positive-time unit steps, namely $X,Y$ in
the diagonal slicing. An increment in a negative direction is
exactly the negative of the corresponding forward increment
ending at $z$. This rewriting is local, even if previous cuts
have been made. Choose a point on the common prefix in the
intact neighborhood, express each increment as a difference of
two complete primitives from that point, cancel their common
initial marked terms, and apply Lemma \ref{cur:deformation}
inside the neighborhood. Thus no earlier path or cut is altered.

Use the following common endpoints for the high stencil:
\begin{itemize}
\item For directions $a,b$, start at $z$ and append two
      after-twist steps in the other direction to end at
      $z+2a+2b$.
\item For directions $-a,-b$, start at $z-2a-2b$, prepend
      two before-twist steps in the other direction, and end
      at $z$.
\item For directions $a,-b$, use the trajectory starting at
      $z-2b$ with steps $b,b,a,a$.
\end{itemize}
The case $-a,b$ is the same with $a,b$ interchanged.
Backward differencing changes the relative sign exactly as
required to obtain the coefficient for the two positive
directions.

Approach the common start by the same prefix, with a final
positive-direction segment. A simple replacement can first bend
around the stencil and approach from the required side, even
when the original prefix arrived from later times. In a box
with spare room, choose lattice corridors preserving the
entrance and reserve a clear rectangle for the last positive
steps. All marker lifts now differ by the short angle between
$a$ and $b$ in their common time frame; no extra $2\pi$ turn
is introduced. The low case is the same construction with one
increment. The explicit prefix charge is $Q^\epsilon$, with
$\epsilon=1$ or $-1$ according to the reflected type.

We next consolidate diagonal chains in the clear band. Keep as
common background every chain except the explicit prefix
steps between the common start and the marker. In particular,
the charge arriving at the common start is included in this
background. Using real charges modulo $2\pi$, route it to carry
a total phase $\zeta$ on a common positive-time seam before
the start, and a phase $\xi$ from the start through the
comparison block and to the future seam. These are the net
oriented charges through ordinary sections, and
\[
 \zeta/\xi=Q^\epsilon.
\]
All endpoints are preserved. Joining and funneling the chains
can be done outside the sandwich cuts but inside the larger
vacant margins, including transverse lattice displacements.
Choose the cuts through level vertices of the positive-time
seam with many positive steps on each side. This construction
does not require a monotone path between the original remote
cut positions.

Inside the comparison rectangle, first route the $\xi$
background along one common skirt on the higher-space-coordinate
side of both positive-step paths, avoiding every possible
marker. In $X,Y$ coordinates it runs around the northwest side
of the small positive rectangle between the common endpoints.
Tile ice conservation then moves this skirt onto either
compared path, always on the same, left spin side of a broken
source edge. At a marker it therefore multiplies $L_-$ or
$L_+$ on a fixed side by the diagonal twist. This multiplication
is one common scalar.

It remains to justify that the change from the original
background routing to this common skirt has a common phase.
Their difference is a closed tile-edge chain in the larger
slab, hence a weighted sum of face boundaries. Tile ice
conservation cancels every face contribution except at a
source edge, whose signed spin jump is fixed and equal to
two relative to its positive step. All prospective source
edges lie in the same winding region of this common
closed-chain difference. Indeed other original chains avoid
the intact neighborhood; the chosen prefix approaches the
common start from the past without entering the positive
rectangle; the replacement before the start follows that
same final incoming line with net charge $\zeta$; and the
common skirt and outgoing seam stay outside the rectangle.
No edge of the routing difference separates the interiors
of two possible marker edges. Extend the coincident clear
incoming and outgoing parts by many lattice edges if needed.
The face coefficient seen by the source is consequently
identical in every stencil term. Summing conservation gives
one common phase. There are no contributions from other
defects, since all changes remain inside the clear band and
its margins. The same reasoning can be phrased as a sequence
of path slides consistently kept on one side of the entire
marker rectangle. Every phase remains unitary.

We have now obtained exactly the before- and after-twist
operators in \eqref{bulk:MN}, with their required relative
phase, the compensating even pairs described above, and
one common bounded scalar. The two exterior states are
independent of which term is used. Theorem
\ref{bulk:angular-estimate} gives the claimed bound.

For later repeated cuts, note also how these states glue.
A consolidated chain at a diagonal cut ends at a shared
level vertex. Reflection conjugates its charge to the
inverse in the mirrored direction, preserving the spin
convention. Joining after reversing path direction therefore
continues with the same charge and creates no endpoint at
the cut. At preliminary tile-axis cuts, route the chains
transversely over the line through its level vertices, with
no segment lying on the line inside the spare margin; the
same argument applies. Reflections preserve the high or
low type at every other intact stencil point. This proves
the reflected assertion and the compatibility needed for
successive extractions.
\end{proof}

The proposition is a local norm estimate. In later applications
the number of available clear slabs and the sizes of their
exterior norms are estimated separately by the positive
partition comparisons; the local extraction itself makes no
uniform bound on those norms.
